\documentclass[10pt,a4paper,landscape]{article}
\usepackage[utf8]{inputenc}
\usepackage[T1]{fontenc}
\usepackage[french]{babel}
\frenchbsetup{StandardLists=true}
\usepackage[landscape,left=1cm,right=1cm,top=1.5cm,bottom=1.5cm]{geometry}
\usepackage{multirow,tabularx,booktabs}
\usepackage[dvipsnames]{xcolor}
\usepackage{fouriernc}
\usepackage{chemmacros}
\chemsetup{modules=all}
\usepackage{fancyhdr}
\usepackage{colortbl}
\usepackage{multicol}
\setlength{\columnsep}{8mm}
\setlength{\columnseprule}{0.4pt}
\pagestyle{fancy}
\renewcommand\headrulewidth{1pt}
\fancyhead[L]{Lycée Denis Diderot}
\fancyhead[R]{Classe de 3PM}
\fancyfoot[C]{}
\fancyfoot[R]{Année 2025-2026}
\fancyfoot[L]{Alexis RUIZ}
\setlength{\parindent}{0pt}
\renewcommand{\arraystretch}{1.35}

% Couleurs
\definecolor{exoheader}{RGB}{30,90,160}
\definecolor{exolight}{RGB}{220,235,255}
\definecolor{totalrow}{RGB}{255,230,200}

% Macro en-tête d'exercice pleine largeur
\newcommand{\exotitle}[2]{%
  \noindent\colorbox{exoheader}{%
    \parbox{\dimexpr\linewidth-2\fboxsep\relax}{%
      \textcolor{white}{\textbf{\large #1 \hfill #2}}%
    }%
  }%
}

\begin{document}

\begin{center}
	{\LARGE\textbf{GRILLE DE CORRECTION}}\\[1mm]
	{\large Chapitre 4 -- Les ions \hspace{1cm} 3PM \hspace{1cm} \textbf{Total : /30 points}}
\end{center}

\vspace{2mm}

\begin{multicols}{2}

% ===================== EXERCICE 1 =====================
\exotitle{Exercice 1 : Connaissances de cours}{/6 pts}

\vspace{1mm}

\begin{tabularx}{\linewidth}{|c|X|c|}
	\hline
	\rowcolor{exolight}
	\textbf{Q} & \textbf{Réponse attendue} & \textbf{Pts} \\
	\hline
	1 & Atome (ou groupe d'atomes) ayant gagné ou perdu des électrons $\Rightarrow$ charge électrique & 1 \\
	\hline
	\multirow{2}{*}{2} & \textbf{Cation} : ion chargé \textbf{positivement}, a \textbf{perdu} des électrons & 1 \\
	\cline{2-3}
	& \textbf{Anion} : ion chargé \textbf{négativement}, a \textbf{gagné} des électrons & 1 \\
	\hline
	3 & En gagnant ou en perdant des électrons (le noyau ne change pas) & 1 \\
	\hline
	4 & Deux ions parmi : \ch{Cl-}, \ch{Cu^{2+}}, \ch{Fe^{3+}}, \ch{Fe^{2+}}, \ch{Al^{3+}} \ldots \quad (0,5 pt chacun) & 1 \\
	\hline
	5 & \textbf{VRAI} -- somme charges(+) = somme charges(--) & 1 \\
	\hline
	\rowcolor{totalrow}
	\multicolumn{2}{|r|}{\textbf{Sous-total Exercice 1}} & \textbf{/6} \\
	\hline
\end{tabularx}

\vspace{3mm}

% ===================== EXERCICE 2 =====================
\exotitle{Exercice 2 : Composition de l'ion aluminium \ch{Al^{3+}}}{/5 pts}

\vspace{1mm}

\begin{tabularx}{\linewidth}{|c|X|c|}
	\hline
	\rowcolor{exolight}
	\textbf{Q} & \textbf{Réponse attendue} & \textbf{Pts} \\
	\hline
	1 & \textbf{Cation} (charge positive +3) \quad 0,5 pt réponse + 0,5 pt justification & 1 \\
	\hline
	2 & \textbf{13 protons} (le nombre de protons ne change pas) & 1 \\
	\hline
	3 & \textbf{10 électrons} \quad Calcul : $13 - 3 = 10$ \quad (0,5 pt calcul + 1 pt réponse) & 1,5 \\
	\hline
	4 & A \textbf{perdu 3 électrons} \quad (1 pt ``perdu'' + 0,5 pt ``3'') & 1,5 \\
	\hline
	\rowcolor{totalrow}
	\multicolumn{2}{|r|}{\textbf{Sous-total Exercice 2}} & \textbf{/5} \\
	\hline
\end{tabularx}

\vspace{3mm}

% ===================== EXERCICE 3 =====================
\exotitle{Exercice 3 : Compléter le tableau}{/5 pts}

\vspace{1mm}

{\footnotesize\textit{1 pt par ligne -- 0,25 pt par case correcte}}

\vspace{1mm}

\begin{tabularx}{\linewidth}{|X|c|c|c|c|c|}
	\hline
	\rowcolor{exolight}
	\textbf{Particule} & \textbf{Protons} & \textbf{Électrons} & \textbf{Atome/Ion} & \textbf{Formule} & \textbf{/1 pt} \\
	\hline
	Calcium        & 20 & 20 & Atome & \ch{Ca}      & 1 \\
	\hline
	Ion calcium    & 20 & 18 & Ion   & \ch{Ca^{2+}} & 1 \\
	\hline
	Oxygène        & 8  & 8  & Atome & \ch{O}       & 1 \\
	\hline
	Ion oxyde      & 8  & 10 & Ion   & \ch{O^{2-}}  & 1 \\
	\hline
	Ion cuivre (II)& 29 & 27 & Ion   & \ch{Cu^{2+}} & 1 \\
	\hline
	\rowcolor{totalrow}
	\multicolumn{5}{|r|}{\textbf{Sous-total Exercice 3}} & \textbf{/5} \\
	\hline
\end{tabularx}

\vspace{3mm}

% ===================== EXERCICE 4 =====================
\exotitle{Exercice 4 : Tests de reconnaissance}{/6 pts}

\vspace{1mm}

\begin{tabularx}{\linewidth}{|c|X|c|}
	\hline
	\rowcolor{exolight}
	\textbf{Q} & \textbf{Réponse attendue} & \textbf{Pts} \\
	\hline
	\multirow{2}{*}{1} & \textbf{Ion fer (III) \ch{Fe^{3+}}} & 1 \\
	\cline{2-3}
	& Précipité rouille/brun-orangé avec NaOH & 1 \\
	\hline
	\multirow{2}{*}{2} & \textbf{Ion chlorure \ch{Cl-}} & 1 \\
	\cline{2-3}
	& Précipité blanc noircissant à la lumière avec \ch{AgNO3} & 1 \\
	\hline
	\multirow{2}{*}{3} & \textbf{Ion cuivre (II) \ch{Cu^{2+}}} & 1 \\
	\cline{2-3}
	& Précipité bleu avec NaOH & 1 \\
	\hline
	\rowcolor{totalrow}
	\multicolumn{2}{|r|}{\textbf{Sous-total Exercice 4}} & \textbf{/6} \\
	\hline
\end{tabularx}

\vspace{3mm}

% ===================== EXERCICE 5 =====================
\exotitle{Exercice 5 : Tableau des tests de reconnaissance}{/4 pts}

\vspace{1mm}

{\footnotesize\textit{1 pt par ligne -- 0,5 pt réactif + 0,5 pt résultat. Accepter ``soude'' pour NaOH.}}

\vspace{1mm}

\begin{tabularx}{\linewidth}{|X|c|X|c|}
	\hline
	\rowcolor{exolight}
	\textbf{Ion} & \textbf{Réactif} & \textbf{Résultat positif} & \textbf{/1 pt} \\
	\hline
	Ion chlorure \ch{Cl-}       & \ch{AgNO3} & Précipité blanc noircissant à la lumière & 1 \\
	\hline
	Ion cuivre (II) \ch{Cu^{2+}}& NaOH       & Précipité bleu & 1 \\
	\hline
	Ion fer (II) \ch{Fe^{2+}}   & NaOH       & Précipité vert & 1 \\
	\hline
	Ion fer (III) \ch{Fe^{3+}}  & NaOH       & Précipité rouille (brun-orangé) & 1 \\
	\hline
	\rowcolor{totalrow}
	\multicolumn{3}{|r|}{\textbf{Sous-total Exercice 5}} & \textbf{/4} \\
	\hline
\end{tabularx}

\vspace{3mm}

% ===================== EXERCICE 6 =====================
\exotitle{Exercice 6 : Analyse de l'ion chlorure \ch{Cl-}}{/4 pts}

\vspace{1mm}

\begin{tabularx}{\linewidth}{|c|X|c|}
	\hline
	\rowcolor{exolight}
	\textbf{Q} & \textbf{Réponse attendue} & \textbf{Pts} \\
	\hline
	\multirow{2}{*}{1} & \textbf{17 protons} (le nombre de protons ne change pas) & 1 \\
	\cline{2-3}
	& \textbf{18 électrons} \quad ($17 + 1$ électron gagné) & 1 \\
	\hline
	\multirow{2}{*}{2} & \textbf{Anion} (charge négative --1) & 1 \\
	\cline{2-3}
	& A gagné 1 électron $\Rightarrow$ charge négative & 1 \\
	\hline
	\rowcolor{totalrow}
	\multicolumn{2}{|r|}{\textbf{Sous-total Exercice 6}} & \textbf{/4} \\
	\hline
\end{tabularx}

\end{multicols}

\end{document}
